% Base16 Tomorrow Light colors
\definecolor{base00}{HTML}{ffffff} % Background
\definecolor{base01}{HTML}{e0e0e0} % Lighter background
\definecolor{base02}{HTML}{d6d6d6} % Selection
\definecolor{base03}{HTML}{8e908c} % Comments
\definecolor{base04}{HTML}{969896} % Dark foreground
\definecolor{base05}{HTML}{4d4d4c} % Default foreground
\definecolor{base06}{HTML}{282a2e} % Light foreground
\definecolor{base07}{HTML}{1d1f21} % Black
\definecolor{base08}{HTML}{c82829} % Red
\definecolor{base09}{HTML}{f5871f} % Orange
\definecolor{base0A}{HTML}{eab700} % Yellow
\definecolor{base0B}{HTML}{718c00} % Green
\definecolor{base0C}{HTML}{3e999f} % Cyan
\definecolor{base0D}{HTML}{4271ae} % Blue
\definecolor{base0E}{HTML}{8959a8} % Purple
\definecolor{base0F}{HTML}{a3685a} % Brown

\lstdefinestyle{base16python}{
  language=Python,
  morekeywords={pass} % doesn't seem to work..
  backgroundcolor=\color{base00},
  basicstyle={\renewcommand\ttdefault{cmtt}\ttfamily\small\color{base05}},
  keywordstyle=\color{base0E},
  commentstyle=\color{base03},
  stringstyle=\color{base0B},
  numberstyle=\scriptsize\color{base03},
  identifierstyle=\color{base05},
  showstringspaces=false,
  breaklines=true,
  postbreak=\space,
  frame=none,
  rulecolor=\color{base01},
  tabsize=4,
  columns=fullflexible,
  keepspaces=true,
  moredelim=[l][\color{base0B}]{@},
}

% Inline code at 95% of the surrounding text, so that it scales in table cells and captions; the size is set
% outside the listing, which applies its basic style more than once
\makeatletter
\newcommand*{\pyinlinesize}{\fontsize{\strip@pt\dimexpr\f@size pt*19/20\relax}{\f@baselineskip}\selectfont}
\makeatother
\newcommand*{\pyinlinestyle}{\renewcommand\ttdefault{cmtt}\ttfamily}
% File and directory names, in the size and font of inline code
\newcommand*{\filename}[1]{{\pyinlinesize\pyinlinestyle #1}}

% A listing opening a top-aligned minipage already sits a line below the preceding text, so its space above
% is replaced by a small negative skip, measured to leave the same gap above as below.
\makeatletter
\lstnewenvironment{python}[1][]{\lstset{style=base16python}\if@minipage\lstset{aboveskip=-3pt}\fi\lstset{#1}}{}
\makeatother

% Listings side by side: each listingcol is a top-aligned minipage, and listingrow separates the columns by
% faint dotted rules spanning the tallest column; listingrow* sets one listing in columns, without rules.
\makeatletter
\newcount\listing@cols
\newdimen\listing@ht
\newdimen\listing@dp
\@for\listing@i:={1,2,3,4}\do{\expandafter\newbox\csname listing@box\listing@i\endcsname}
\newcommand*{\listing@box}[1]{\csname listing@box#1\endcsname}
\newif\iflisting@rule
\newenvironment{listingrow}{\par\global\listing@cols=0\relax\listing@ruletrue}{\listing@typeset}
\newenvironment{listingrow*}{\par\global\listing@cols=0\relax\listing@rulefalse}{\listing@typeset}
\newcommand*{\listing@typeset}{%
  \listing@ht=0pt\relax \listing@dp=0pt\relax
  \count@=0\relax
  \loop\ifnum\count@<\listing@cols \advance\count@ 1\relax
    \ifdim\ht\listing@box{\the\count@}>\listing@ht \listing@ht=\ht\listing@box{\the\count@}\fi
    \ifdim\dp\listing@box{\the\count@}>\listing@dp \listing@dp=\dp\listing@box{\the\count@}\fi
  \repeat
  \noindent\count@=0\relax
  \loop\ifnum\count@<\listing@cols \advance\count@ 1\relax
    \ifnum\count@>1 \hfill\iflisting@rule\listing@rule\hfill\fi\fi
    \box\listing@box{\the\count@}%
  \repeat
  \par\medskip}
\newcommand*{\listing@rule}{%
  \tikz[baseline=0pt]\draw[base03, densely dotted, line width=0.5pt]
    (0pt,\the\listing@ht) -- (0pt,-\the\listing@dp);}
\newenvironment{listingcol}[1]
  {\global\advance\listing@cols 1\relax
   \global\setbox\listing@box{\the\listing@cols}=\hbox\bgroup\begin{minipage}[t]{#1}}
  {\end{minipage}\egroup}
\makeatother

% Colour of inline code: base05, or the markup colour inside change markup (tex/change-tracking.tex)
\providecommand*{\codecolor}{base05}

% for use in inferrule
\newcommand{\pytext}[1]{%
  {\renewcommand\ttdefault{cmtt}\ttfamily\small\color{\codecolor}#1}%
}
\newcommand{\pyinlinetext}[2]{\text{{\renewcommand\ttdefault{cmtt}\ttfamily\upshape#1\color{\codecolor}#2}}}
% \pymath: math-safe inline code (no \lstinline, so it works inside an
% inferrule). Unlike \pyinline it does not read its argument verbatim, so the
% caller escapes special characters, e.g. \pymath{\_\_name\_\_}.
\newcommand*{\pymath}[1]{%
  \mathchoice
    {\pyinlinetext{}{#1}}%
    {\pyinlinetext{}{#1}}%
    {\pyinlinetext{\scriptsize}{#1}}%
    {\pyinlinetext{\tiny}{#1}}%
}
\newcommand{\pyinline}[1]{%
  \ifmmode \pymath{#1}\else
    \ifcodeinmarkup
      {\pyinlinesize\mbox{\lstinline[style=base16python,basicstyle={\pyinlinestyle\color{\codecolor}},keywordstyle=\color{\codecolor},stringstyle=\color{\codecolor},identifierstyle=\color{\codecolor}]{#1}}}%
    \else
      {\pyinlinesize\mbox{\lstinline[style=base16python,basicstyle={\pyinlinestyle\color{base05}}]{#1}}}%
    \fi
  \fi
}
